\section{Framework}
\label{sec:framework}

% ---------------------------------------------------------------
\subsection{Contact order parameter $\rcontact$}
\label{sec:rho}

A robotic manipulation episode is, at root, a sequence of contact
events: the robot approaches a passive object, applies enough force
to support it against gravity, transports it, and releases it onto
another passive support. The full pipeline of perception, planning,
and control is organized around a single graded property of the
robot--object contact at each instant: \emph{the ratio of the
normal contact force the manipulator applies to the normal force
required for the object to be in static equilibrium}. We call
this dimensionless ratio the \emph{contact order parameter}
$\rcontact$ and take it to be the central interface variable of
the paper. Subsequent subsections
build the spectral, geometric, and control structure of the
five-phase cycle on top of $\rcontact$; the present subsection
defines it.

\begin{definition}[Contact order parameter]
\label{def:rho}
For a body $b$ in instantaneous contact with the environment via a
contact set with outward unit normal $\hat{\mathbf n}$, let
$f_{n} \ge 0$ denote the magnitude of the contact force projected
onto $\hat{\mathbf n}$, and let $f_{n}^{\star} > 0$ denote the
normal force required for $b$ to be in static equilibrium under
gravity and any task-relevant external load. The
\emph{contact order parameter} of $b$ is
\begin{equation}
\label{eq:rho-def}
\rcontact \;\coloneqq\; \frac{f_{n}}{f_{n}^{\star}} \;\in\; [0, \infty).
\end{equation}
\end{definition}

The order parameter $\rcontact$ is dimensionless and non-negative.
We adopt the following sign convention throughout the paper:
\begin{itemize}[leftmargin=*]
  \item $\rcontact < 1$ --- the contact force is insufficient to
    support $b$; the body is slipping or out of contact.
  \item $\rcontact = 1$ --- the static-equilibrium threshold; this
    level set defines the \emph{switching surface}
    $\swsurf = \{\rcontact = 1\}$, formalized in \cref{sec:Sigma}.
  \item $\rcontact > 1$ --- the contact force exceeds the
    equilibrium requirement; the manipulator is supporting the
    body against its load.
\end{itemize}

Two literature anchors place $\rcontact$ in context. In nonsmooth
contact mechanics one classically tracks a discrete contact-mode
indicator that toggles between \emph{contact} and \emph{no
contact} with complementarity conditions on the force and the gap
function~\citep{Brogliato2016}. The order parameter $\rcontact$
generalizes that indicator by replacing the binary toggle with a
continuous measure that smooths the transition through
$\rcontact = 1$. In grasping and force-closure analysis one
analogously tracks the friction-cone margin between the applied
contact wrench and the boundary of the friction
cone~\citep{MurrayLiSastry1994}; the regime $\rcontact > 1$
corresponds to wrenches strictly inside that
margin, while $\rcontact = 1$ marks its boundary. Both
characterizations are recovered by appropriate choice of the
denominator $f_{n}^{\star}$.

\begin{example}[Parallel-jaw grasp of a rigid box]
\label{ex:parallel-jaw}
Consider a parallel-jaw gripper grasping a rigid box of mass $m$
in a uniform gravitational field of magnitude $g$. Each finger
applies a normal force of magnitude $f_{n}$ to a vertical face of
the box; let $\mu$ denote the static friction coefficient between
finger and box. The static equilibrium condition for the box is
that the total tangential friction available at the two finger
contacts supports the weight,
\[
2 \mu f_{n} \;\ge\; m g,
\]
with equality at the slipping threshold. The grasp-specific
order parameter is therefore
\begin{equation}
\label{eq:rho-grasp}
\rcontact_{\mathrm{grasp}} \;\coloneqq\; \frac{2 \mu f_{n}}{m g},
\end{equation}
which realizes \cref{def:rho} with $f_{n}^{\star} = m g / (2 \mu)$.
The three regimes specialize as expected:
$\rcontact_{\mathrm{grasp}} < 1$ corresponds to the box slipping
through the fingers, $\rcontact_{\mathrm{grasp}} = 1$ to the
critical grasp force, and $\rcontact_{\mathrm{grasp}} > 1$ to a
grasp that supports the box and can transmit acceleration to it.
\end{example}

The remainder of \cref{sec:framework} develops three structural
consequences of \cref{def:rho}: the spectral content of multi-body
contact (\cref{sec:specgap}, where $\specgap$ characterizes the
connectivity of the contact graph), the geometry of crossings of
the level set $\rcontact = 1$ (\cref{sec:Sigma}, the switching
surface), and the control implications of those crossings
(\cref{sec:three-term} onward).

% ---------------------------------------------------------------
\subsection{Spectral gap $\specgap$ of the contact graph}
\label{sec:specgap}

The order parameter $\rcontact$ of \cref{def:rho} is defined at a
single contact between one body and a counterpart support.
Manipulation tasks of practical interest
typically involve more than one such contact: a stack of two or
more blocks transmits load through a chain of body--body
contacts, a grasped object simultaneously contacts two fingers
and (transiently) a placement surface, and an in-hand
manipulation maneuver relies on several finger--object contacts
arranged around the object. In each case the relevant question
is no longer the value of $\rcontact$ at any individual contact
alone but the connectivity of the entire \emph{contact graph}
that links
the bodies in play to the supporting environment. The spectral
gap $\specgap$ measures this connectivity.

\paragraph{Contact graph.}
At a fixed instant $t$, let
\begin{equation}
\label{eq:contact-graph}
  G_{c}(t) \;=\; (V, E, w),
\end{equation}
where the vertex set $V$ collects the bodies currently in
contact with at least one other body or with the environment,
the edge set $E \subseteq \binom{V}{2}$ collects pairs of
bodies in direct contact, and the edge weight
$w_{ij} \ge 0$ on $\{i, j\} \in E$ is the contact stiffness at
the $(i, j)$ contact, equivalently the magnitude of the contact
force at that contact. (These two interpretations are
equivalent up to a positive scaling: under the linearized
contact model used throughout \cref{sec:framework} a contact at
$\rcontact > 1$ transmits force in proportion to local stiffness,
so taking either as the edge weight produces the
same connectivity structure.) We adopt the convention that the
environment (a fixed inertial support such as a tabletop)
appears as a distinguished vertex in $V$.

The weighted Laplacian of $G_{c}$ is the standard
spectral-graph-theoretic object
\begin{equation}
\label{eq:laplacian-def}
  L_{c} \;\coloneqq\; D - W,
\end{equation}
where $W \in \mathbb{R}^{|V| \times |V|}$ is the weighted
adjacency matrix with entries $W_{ij} = w_{ij}$ for $\{i, j\}
\in E$ and zero otherwise, and $D$ is the diagonal degree
matrix with $D_{ii} = \sum_{j} W_{ij}$. The matrix $L_{c}$ is
symmetric positive semidefinite and admits the quadratic-form
representation
\begin{equation}
\label{eq:laplacian-quadform}
  \mathbf{x}^{\top} L_{c}\, \mathbf{x}
  \;=\; \tfrac{1}{2}\!\sum_{\{i, j\} \in E}\! w_{ij}\,(x_{i} - x_{j})^{2},
  \qquad \mathbf{x} \in \mathbb{R}^{|V|},
\end{equation}
from which the all-ones vector $\mathbf{1}_{V}$ visibly lies in
the kernel of $L_{c}$.

\paragraph{Spectral gap.}
The eigenvalues of $L_{c}$ are non-negative and may be ordered
$0 = \mu_{0} \le \mu_{1} \le \dots \le \mu_{|V|-1}$. We denote
the second-smallest eigenvalue by
\begin{equation}
\label{eq:specgap-def}
  \specgap(L_{c}) \;\coloneqq\; \mu_{1},
\end{equation}
in keeping with the macro definition of \cref{sec:notation}.
Following \citet{Fiedler1973}, $\specgap$ is the
\emph{algebraic connectivity} of $G_{c}$; the broader spectral
graph-theoretic background is collected in \citet{Chung1997}.
The spectral gap is the central quantity that promotes
$\rcontact$ from a single-contact scalar to a multi-body
connectivity measure: it asks whether every body in $V$ is tied,
possibly through a chain of intermediate contacts, to the
supporting environment.

\begin{lemma}[Spectral gap interpretation]\label{lem:specgap-interp}
Let $L_{c}$ be the weighted Laplacian of the contact graph
$G_{c} = (V, E, w)$ with non-negative edge weights. Then
$\specgap(L_{c}) = 0$ if and only if $G_{c}$ is disconnected,
i.e., at least one body in $V$ is not connected via a chain of
positively-weighted contacts to the supporting environment.
Equivalently, $\specgap(L_{c}) > 0$ if and only if every body
in $V$ belongs to a single connected support structure.
\end{lemma}

The proof, a standard consequence of the spectral theory of
weighted graphs, is deferred to
\cref{app:proof:specgap-interp}.

\paragraph{Two-block stack.}
The simplest nontrivial example exhibits all three elements
($\rcontact$, $L_{c}$, and $\specgap$) at once. Consider two
cubes $b_{1}$, $b_{2}$ stacked vertically with $b_{1}$ resting
on a table and $b_{2}$ resting on $b_{1}$. The vertex set is
$V = \{b_{1}, b_{2}, \text{table}\}$ and the edge set is
$E = \{\{b_{2}, b_{1}\},\,\{b_{1}, \text{table}\}\}$ with edge
weights $w_{12}$ and $w_{1t}$, respectively. The associated
Laplacian (ordering vertices $b_{1}, b_{2}, \text{table}$) is
\begin{equation}
\label{eq:two-block-L}
  L_{c} \;=\;
  \begin{pmatrix}
    w_{12} + w_{1t} & -w_{12} & -w_{1t} \\
    -w_{12} & w_{12} & 0 \\
    -w_{1t} & 0 & w_{1t}
  \end{pmatrix}.
\end{equation}
Direct computation shows that $L_{c}$ has eigenvalue $0$ with
eigenvector $\mathbf{1}_{V}$, plus two further eigenvalues
whose product equals
$\det(L_{c}|_{\mathbf{1}^{\perp}}) = w_{12}\,w_{1t} +
w_{12}\,w_{1t}$ up to scaling, and which are both positive
exactly when both $w_{12}$ and $w_{1t}$ are positive. The
spectral gap therefore satisfies
\begin{equation}
\label{eq:two-block-specgap}
  \specgap(L_{c}) > 0
  \quad\Longleftrightarrow\quad
  w_{12} > 0 \;\;\text{and}\;\; w_{1t} > 0.
\end{equation}
Translated back into the order-parameter language of
\cref{def:rho}, $\specgap > 0$ requires that both
$\rcontact_{b_{2} \to b_{1}} \ge 1$ (the upper cube does not
slip on the lower) and $\rcontact_{b_{1} \to \text{table}} \ge
1$ (the lower cube does not slip on the table). The spectral
gap is positive precisely when every contact in the stack
satisfies $\rcontact \ge 1$, so that the chain forms a single
connected support structure rooted in the environment;
it collapses to zero the moment any single contact in the chain
fails. This is the multi-body extension of $\rcontact$ promised
above and the connectivity criterion that the controller will
maintain through phase 3 of the cycle (\cref{sec:three-term}).

% ---------------------------------------------------------------
\subsection{Switching surface $\swsurf$ and the two-crossing lemma}
\label{sec:Sigma}

The level set
\begin{equation}
\label{eq:Sigma-def}
  \swsurf \;\coloneqq\; \{\rcontact = 1\}
\end{equation}
introduced in passing in \cref{sec:rho} plays the role of a
\emph{switching surface} in the joint state space spanned by
the order parameter $\rcontact$ together with the body
configuration $x$ and velocity $\dot x$. We treat
$\swsurf$ as the codimension-one subset
$\{(\rcontact, x, \dot x) : \rcontact = 1\}$, although, since
none of the qualitative structure of this subsection depends on
the additional coordinates $(x, \dot x)$, it is sometimes
sufficient to view $\swsurf$ as a level set of $\rcontact$ on
its own. The complement of $\swsurf$ partitions the joint state
space into the region $\{\rcontact < 1\}$ (no contact, or contact
without sufficient force to support the object) on one side and
the region $\{\rcontact > 1\}$ (the manipulator supports the
object) on the other. The robot's task is to traverse a
trajectory whose
$\rcontact$-projection visits both sides exactly when needed.

\paragraph{Boundary conditions of a pick--place task.}
A pick--place task is specified by an initial pose at time $t =
0$ and a final pose at time $t = T$. At both endpoints the
object rests on a passive support (the initial table at $t = 0$
and a possibly different placement surface at $t = T$); the
robot is not in contact with the object. From the perspective
of the robot--object contact whose order parameter is
$\rcontact$, the boundary conditions are therefore
\begin{equation}
\label{eq:pick-place-BCs}
  \rcontact(0) \;<\; 1, \qquad \rcontact(T) \;<\; 1.
\end{equation}
For the task to actually transport the object from its initial
support to its final support, the robot must, during some
interval of $(0, T)$, support the object against gravity
through its own contact: there must exist some $t^{*} \in (0,
T)$ at which $\rcontact(t^{*}) > 1$. Equation
\eqref{eq:pick-place-BCs} together with this in-task condition
constitutes the standing assumption of this subsection.

\begin{lemma}[Two-crossing lemma]\label{lem:two-crossing}
Let $\rcontact \colon [0, T] \to [0, \infty)$ be the
robot--object contact order parameter along any pick--place
trajectory satisfying the boundary conditions $\rcontact(0) <
1$ and $\rcontact(T) < 1$, with $\rcontact(t) > 1$ on at least
one open subinterval of $(0, T)$. If $\rcontact$ is continuous
on $[0, T]$, then $\rcontact$ crosses the level set $\swsurf =
\{\rcontact = 1\}$ at least twice in $(0, T)$, and any
pick--place trajectory of minimal complexity has exactly two
crossings: one upward, at the contact-establishment time
$t_{1}$, and one downward, at the release time $t_{4}$.
\end{lemma}

The proof, an application of the intermediate value theorem,
is deferred to \cref{app:proof:two-crossing}. The
labels $t_{1}$ and $t_{4}$ are deliberately chosen to align
with the phase numbering of \cref{sec:phase-cycle}: phase 2 is
the bang arc concentrated at $t_{1}$ and phase 4 is the bang
arc concentrated at $t_{4}$.

The two crossings partition $[0, T]$ into three intervals:
\emph{pre-contact} $[0, t_{1}]$ on which $\rcontact < 1$,
\emph{in-contact} $[t_{1}, t_{4}]$ on which $\rcontact \ge 1$,
and \emph{post-contact} $[t_{4}, T]$ on which $\rcontact < 1$
again. Section~\ref{sec:phase-cycle} will further subdivide
the in-contact interval $[t_{1}, t_{4}]$ into the bang arcs
(phases 2 and 4) localized at the two crossings and the
singular arc (phase 3) on the open interior $(t_{1}, t_{4})$
where the spectral-gap maintenance condition $\specgap \ge
\varepsilon$ of \cref{sec:specgap} is in force. The pre- and
post-contact intervals together with the in-contact transport
arc account for the five-phase decomposition of the cycle
introduced in the next section.

The two-crossing lemma should be read as a forced
consequence of two facts only: the continuity of $\rcontact$
along admissible trajectories and the boundary conditions
\eqref{eq:pick-place-BCs} of a pick--place task. No
implementation choice, no assumption on the robot's actuators,
and no assumption on the controller architecture enters the
statement. Continuity of $\rcontact$ is a mild hypothesis: it
holds whenever the contact force evolves continuously in time,
which is the case for any compliant or finite-stiffness contact
model and for any trajectory generated by a controller with
bounded torque rates. The intermediate value theorem
\citep{Rudin1976} then forces the two crossings.

% ---------------------------------------------------------------
\subsection{Three-term control decomposition}
\label{sec:three-term}

The combinatorial picture of \cref{sec:Sigma}---a pick--place
trajectory crosses $\swsurf$ exactly twice, at times $t_{1}$
and $t_{4}$---has a precise control-theoretic counterpart. The
optimal joint torque $u$ along such a trajectory factors into
three structurally distinct contributions: a configuration-only
gravity term, a singular-arc least-action term that lives on
the open intervals away from the crossings, and an impulsive
spectral kick concentrated on $\swsurf$ itself. This subsection
states the decomposition; the proof is deferred to
\cref{app:proof:three-term}.

\paragraph{Optimal control problem.}
Throughout this subsection we consider a pick--place optimal
control problem (OCP) on the time interval $[0, T]$ with state
$x = (\jointvec, \dot{\jointvec}, x_{\mathrm{obj}})$ comprising
the manipulator joint configuration $\jointvec$ and velocity
$\dot{\jointvec}$ together with the object pose
$x_{\mathrm{obj}} \in SE(3)$. Let $L(x, \dot x, u, t)$ be a
running cost (energy, smoothness) and let $\Psi(x_{T})$ be a
terminal cost (placement accuracy). The OCP is
\begin{equation}
\label{eq:ocp}
  \min_{u} \;
  J[u]
  \;=\;
  \int_{0}^{T} L(x, \dot x, u, t)\,dt
  \;+\; \Psi(x_{T}),
\end{equation}
subject to the manipulator--object dynamics $\dot x = f(x, u)$
and to the spectral-gap constraint
\begin{equation}
\label{eq:specgap-constraint}
  \rcontact(t) > 1
  \;\Longrightarrow\;
  \specgap\bigl(L_{c}(t)\bigr) \;\ge\; \varepsilon
\end{equation}
on the in-contact interval, with $\varepsilon$ the spectral-gap
threshold of the notation table (\cref{sec:notation}). The
constraint \eqref{eq:specgap-constraint} encodes the spectral
connectivity requirement of \cref{sec:specgap}: while the robot
supports the object, the contact graph must remain connected with
quantitative margin.

\begin{definition}[Three-term decomposition]\label{def:three-term}
For a pick--place OCP as above, an admissible control input
$u \colon [0, T] \to \mathbb{R}^{n}$ is said to admit a
\emph{three-term decomposition} if it can be written as
\begin{equation}
\label{eq:three-term}
  u(t) \;=\; u_{g}(t) \;+\; u_{a}(t) \;+\; u_{s}(t),
\end{equation}
where:
\begin{enumerate}[label=(\roman*),nosep]
  \item $u_{g}(t)$ is the \emph{gravity-compensation term}: the
    joint torque required to support the manipulator and any
    in-contact object against gravity at the current configuration
    $\jointvec(t)$.
  \item $u_{a}(t)$ is the \emph{least-action term}: the residual
    control effort allocated by the OCP first-order conditions on
    the singular interior of each phase, away from $\swsurf$.
  \item $u_{s}(t)$ is the \emph{spectral-kick term}: an impulsive
    or near-impulsive contribution supported on the
    $\swsurf$-crossing times $\{t_{1}, t_{4}\}$ that establishes
    or releases the contact graph spectral gap.
\end{enumerate}
\end{definition}

\begin{theorem}[Three-term optimal decomposition]\label{thm:three-term}
Under standard regularity assumptions on $L$, $\Psi$, and the
dynamics $f$, the optimal control of the pick--place OCP
\eqref{eq:ocp} admits a three-term decomposition in the sense
of \cref{def:three-term}, with $u_{s}$ supported on
$\{t : x(t) \in \swsurf\}$ and $u_{a}$ satisfying the
Pontryagin necessary conditions on the complement.
\end{theorem}

The proof, an application of Pontryagin's maximum principle
together with the bang--singular structure of OCPs with a
switching constraint, is given in
\cref{app:proof:three-term}.

\paragraph{Per-phase reading.}
Aligned with the labels $t_{1}, t_{4}$ of the two-crossing
lemma, the three terms light up in different combinations
across the five phases that \cref{sec:phase-cycle} will
formalize. On phases 1 (pre-contact, $t \in [0, t_{1}]$) and 5
(post-contact, $t \in [t_{4}, T]$) the trajectory is a pure
singular arc: there is no $\swsurf$-crossing in its interior,
so $u_{s} = 0$ and the control consists of $u_{g}$ plus the
least-action $u_{a}$ alone. On phases 2 (load, at $t = t_{1}$)
and 4 (unload, at $t = t_{4}$) the trajectory is a bang arc
concentrated at the corresponding $\swsurf$-crossing and the
spectral-kick term $u_{s}$ is active; the gravity and
least-action terms remain present but are dominated by the
impulsive contribution. On phase 3 (transport, $t \in (t_{1},
t_{4})$) the trajectory is again a singular arc (so $u_{s} =
0$ on the open interior) but now subject to the additional
spectral-maintenance constraint $\specgap(t) \ge \varepsilon$,
which the inner level of the bi-level optimization in
\cref{sec:bilevel} will enforce. The decomposition therefore
serves both as a qualitative reading of the optimal cycle and
as the template against which the controller of
\cref{sec:bilevel} is structured.

The maximum principle invoked in
\cref{thm:three-term} is the version of
\citet{Pontryagin1962}, and the bang--singular structure at the
$\swsurf$-crossings is the standard consequence of a switching
constraint with discontinuous right-hand side at the level set,
in the sense of \citet{BrysonHo1975}.

% ---------------------------------------------------------------
\subsection{Task-space / null-space decomposition}
\label{sec:task-null}

The three-term decomposition of \cref{sec:three-term} factorizes the
optimal joint torque \emph{temporally}: a configuration-only gravity
term, a singular least-action term on the open intervals, and an
impulsive spectral kick on $\swsurf$. The present subsection
introduces a complementary \emph{spatial} factorization in joint
space: at each instant, the joint velocity splits into a part that
moves the end-effector and a part that rearranges the manipulator
internally without affecting the end-effector. Together with the
temporal three-term reading, this spatial split sets up the
five-phase cycle of \cref{sec:phase-cycle} and the hierarchical
controller of \cref{sec:hierarchical}.

\paragraph{Setup.}
Let the manipulator have $n$ joints with configuration
$\jointvec \in \mathbb{R}^{n}$; for the Franka Emika Panda used in
the experimental staircase of \cref{sec:experiments}, $n = 7$. The
end-effector pose lives on a task manifold of dimension
$m \le n$, typically $m = 6$ for full $SE(3)$ pose tracking. The
\emph{task Jacobian} $J(\jointvec) \in \mathbb{R}^{m \times n}$
relates joint velocity to end-effector velocity:
\begin{equation}
\label{eq:task-jacobian}
  \dot x \;=\; J(\jointvec)\, \dot{\jointvec}.
\end{equation}
Throughout this subsection we assume that $J$ has full row rank $m$
generically along the trajectory; isolated kinematic singularities,
where the rank drops, are handled by standard damped-least-squares
regularization and are not treated separately here.

\paragraph{Decomposition.}
Under the redundancy assumption $n > m$ and full row rank of $J$,
the joint velocity admits the operational-space decomposition
\begin{equation}
\label{eq:qdot-split}
  \dot{\jointvec}
  \;=\; J^{+}(\jointvec)\, \dot x
  \;+\; \bigl(I - J^{+}(\jointvec)\, J(\jointvec)\bigr)\, \eta,
\end{equation}
where $J^{+} \in \mathbb{R}^{n \times m}$ is the Moore--Penrose
pseudo-inverse of $J$ and $\eta \in \mathbb{R}^{n}$ parameterizes
\emph{self-motion}: joint rearrangements that produce zero
end-effector velocity. Equation \eqref{eq:qdot-split} is the
operational-space formulation of \citet{Khatib1987} and the
redundancy-resolution split of \citet{Nakamura1991}; the two
projectors $J^{+} J$ and $I - J^{+} J$ are orthogonal complements
in $\mathbb{R}^{n}$ and exhaust every joint-velocity direction.

\begin{definition}[Task-space and null-space subspaces]\label{def:TS-NS}
The \emph{task-space subspace}, denoted $\TS \subseteq \mathbb{R}^{n}$,
is the range of $J^{\top}(\jointvec)$:
\begin{equation}
\label{eq:TS-def}
  \TS \;=\; \mathrm{range}\!\bigl(J^{\top}(\jointvec)\bigr).
\end{equation}
Joint torques in $\TS$ produce non-zero generalized force at the
end-effector. The \emph{null-space subspace}, denoted
$\NS \subseteq \mathbb{R}^{n}$, is the kernel of $J(\jointvec)$:
\begin{equation}
\label{eq:NS-def}
  \NS \;=\; \ker\!\bigl(J(\jointvec)\bigr).
\end{equation}
Joint velocities in $\NS$ produce zero end-effector velocity.
\end{definition}

\paragraph{Dimensions and orthogonality.}
For full-rank $J$, the rank--nullity theorem and the standard
properties of the Moore--Penrose pseudo-inverse give
\begin{equation}
\label{eq:TS-NS-dims}
  \dim \TS \;=\; m,
  \qquad
  \dim \NS \;=\; n - m,
  \qquad
  \mathbb{R}^{n} \;=\; \TS \oplus \NS,
\end{equation}
with the direct sum orthogonal under the standard inner product.
The two projectors $P_{\TS} = J^{\top}(J J^{\top})^{-1} J$ and
$P_{\NS} = I - P_{\TS}$ are mutually orthogonal idempotents and
together resolve the identity on $\mathbb{R}^{n}$.

\begin{example}[Franka Panda elbow self-motion]\label{ex:franka-elbow}
For the Franka Emika Panda, $n = 7$ joints and a task manifold of
dimension $m = 6$ (full $SE(3)$ end-effector pose) give
$\dim \TS = 6$ and $\dim \NS = 1$. The single null-space direction
is the elbow swivel about the line joining the shoulder and the
wrist: the elbow rotates around this line while the end-effector
pose remains fixed. This one-dimensional self-motion is the only
internal reconfiguration available to the arm at a typical
configuration and is the elementary instance of the decomposition
\eqref{eq:qdot-split} that recurs throughout the paper.
\end{example}

\paragraph{Translation note.}
In the operational-space literature this split is sometimes labeled
differently; we adopt the task-space / null-space terminology of
\citet{Khatib1987} and \citet{Nakamura1991} as the most direct, and
use the macros $\TS$ and $\NS$ consistently throughout the
paper. \Cref{sec:hierarchical} will specialize
\cref{def:TS-NS} per phase of the cycle---identifying, for each
$\phaseop{i}$, what the task variable is and what the null-space
freedom encodes---and will define a commutator on $(\TS, \NS)$
that quantifies dexterity in the sense of \citet{Yoshikawa1985}.
The present subsection states only the decomposition itself;
per-phase content and the dexterity invariant are deferred to
\cref{sec:hierarchical}.

% ---------------------------------------------------------------
\subsection{Bi-level optimization scaffold}
\label{sec:bilevel}

The three-term decomposition of \cref{sec:three-term} and the
task/null split of \cref{sec:task-null} characterize the optimal
control structurally---what shape it has and how it factors. To
move from structural characterization to a computable formulation
requires an explicit optimization scaffold. Manipulation OCPs of
the form \eqref{eq:ocp} naturally split into two coupled levels:
an \emph{outer} trajectory-optimization level over the object's
path through space and time, and an \emph{inner} force-closure
level over contact forces, instantaneously. The present subsection
states this scaffold formally; \cref{sec:hierarchical} specializes
it to the pick--place cycle and \cref{sec:experiments} ties it to
the four experimental stages.

\paragraph{Bi-level optimization preliminaries.}
A bilevel program in the standard sense is an optimization problem
whose constraints include the optimality conditions of a second,
nested optimization problem; an upper-level decision variable
parameterizes the lower-level program, and the lower-level
optimum feeds back into the upper-level objective and constraints.
The reader is referred to the survey of \citet{Colson2007} for a
contemporary overview and to standard texts on nonsmooth contact
mechanics \citep{Brogliato2016} and force-closure analysis
\citep{MurrayLiSastry1994} for the manipulation-specific
instantiation that follows.

\begin{definition}[Bi-level pick--place scaffold]\label{def:bilevel}
The pick--place OCP \eqref{eq:ocp} decomposes into two coupled
optimization levels:
\begin{enumerate}[label=(\roman*),nosep]
  \item \textbf{Outer (trajectory).} Choose a velocity field
    $\veltheta(x, t)$ that transports the object from its
    initial pose $g^{A}$ to its final pose $g^{B} \in SE(3)$,
    minimizing the trajectory cost
    \begin{equation}
    \label{eq:bilevel-outer}
      \min_{\veltheta} \;
      \mathbb{E}\!\left[
        \int_{0}^{T} L(x, \dot x)\,dt
        \;+\; \Psi(x_{T})
      \right]
    \end{equation}
    subject to the continuity equation
    $\partial_{t}\mu_{t} + \nabla\!\cdot\!(\veltheta\, \mu_{t}) = 0$
    governing the evolution of the state distribution $\mu_{t}$,
    and to the spectral-gap maintenance condition
    $\specgap(L_{c}(t)) \ge \varepsilon$ on the in-contact
    interval $[t_{1}, t_{4}]$ identified in \cref{lem:two-crossing}.
  \item \textbf{Inner (force closure).} At every instant
    $t \in [t_{1}, t_{4}]$, choose finger contact forces
    $f(t) = (f_{n}(t), f_{t}(t))$ that maximize the contact
    graph spectral gap
    \begin{equation}
    \label{eq:bilevel-inner}
      \max_{f(t)} \;
      \specgap\!\bigl(L_{c}(f(t))\bigr)
    \end{equation}
    subject to the friction-cone constraint
    $|f_{t}| \le \mu f_{n}$ and the Newton balance
    \begin{equation}
    \label{eq:newton-balance}
      \sum_{i} f_{i}
      \;=\;
      m\bigl(\ddot{x}_{\mathrm{des}}(t) - g\bigr)
      \;+\; \Fwind(t),
    \end{equation}
    where $\ddot{x}_{\mathrm{des}}(t)$ is the desired object
    acceleration prescribed by the outer level, $m$ is the object
    mass, $g$ the gravitational acceleration, and $\Fwind(t)$ the
    adversarial wind force of \cref{sec:curriculum}.
\end{enumerate}
\end{definition}

\begin{proposition}[Bi-level coupling]\label{prop:bilevel-coupling}
In \cref{def:bilevel}, the outer and inner levels are coupled
bidirectionally:
\begin{enumerate}[label=(\roman*),nosep]
  \item \emph{Outer to inner.} The desired object acceleration
    $\ddot{x}_{\mathrm{des}}(t)$ enters the inner level through
    the Newton balance \eqref{eq:newton-balance}: faster outer
    trajectories demand larger contact forces and consume
    friction-cone margin.
  \item \emph{Inner to outer.} Force-closure feasibility imposes
    a state-dependent acceleration bound
    \begin{equation}
    \label{eq:accel-bound}
      \|\ddot{x}\|
      \;\le\;
      \alpha\bigl(\jointvec, f_{n,\max}\bigr),
    \end{equation}
    which constrains the outer optimization. The function
    $\alpha$ depends on the current configuration $\jointvec$ and
    the normal-force budget $f_{n,\max}$ available at the contact;
    its explicit form for the parallel-jaw grasp under wind is
    derived in \cref{sec:hierarchical}.
\end{enumerate}
\end{proposition}

In practice the coupling of \cref{prop:bilevel-coupling} is
resolved either by alternating updates---an outer trajectory step,
followed by an inner force-closure step at each instant of the
proposed trajectory, followed by a corrective outer step---or by a
single-shot reformulation when the inner program admits a
closed-form maximum-spectral-gap solution. The choice of solver
is deferred to \cref{sec:hierarchical}; the present subsection
fixes only the structural decomposition.

\paragraph{Friction model.}
The friction-cone constraint $|f_{t}| \le \mu f_{n}$ that appears
in the inner program is the convex-cone form of the Coulomb
friction law. In the unified set-valued formulation of
\citet{Stewart2000}---also adopted in \citet{Brogliato2016}---the
friction force is a subdifferential inclusion
\begin{equation}
\label{eq:friction-di}
  f_{t}(t) \;\in\; -\mu\, f_{n}(t)\,
  \partial \!\left\| v_{t}(t) \right\|,
\end{equation}
where $v_{t}$ is the tangential relative velocity at the contact
and $\partial\|\cdot\|$ is the subdifferential of the Euclidean
norm. When $v_{t} \neq 0$ (slip) this inclusion selects a unique
friction force opposing motion; when $v_{t} = 0$ (stick) it yields
the full friction disc $\|f_{t}\| \le \mu f_{n}$ as a set-valued
static-friction reaction. The friction cone $\frictioncone$
of the notation table \eqref{tab:notation} is the closure of this
set-valued condition: $f \in \frictioncone$ iff $|f_{t}| \le
\mu f_{n}$ and $f_{n} \ge 0$.

\paragraph{Forward link to grasping.}
Within the bi-level scaffold of \cref{def:bilevel}, grasping is the
deliberate creation of a contact singularity at the boundary of
the inner level. Phase 2 (load) drives the inner solution from no
contact ($\specgap = 0$) up to the boundary $\specgap = \varepsilon$
of the feasible region for the inner spectral-gap constraint,
while phase 4 (unload) drives it back
across the same boundary in the opposite direction. The two
boundary crossings of the inner program are exactly the two bang
arcs identified in \cref{def:three-term}: $u_{s}$ at $t_{1}$ is
the impulsive control input that induces the upward boundary
crossing of the inner spectral-gap program, and $u_{s}$ at
$t_{4}$ induces the downward crossing. Phase 3 then operates the
inner program on its singular interior, maintaining
$\specgap(t) \ge \varepsilon$ for all $t \in (t_{1}, t_{4})$.
This identification of the two-crossing geometry of $\swsurf$
with the singular-boundary geometry of the inner program is the
bridge from the structural results of \cref{sec:Sigma} and
\cref{sec:three-term} to the controller architecture of
\cref{sec:hierarchical}.

% ---------------------------------------------------------------
\subsection{Smooth-contact regime and the contact-implicit limit}
\label{sec:smooth-contact}

The framework of \cref{sec:rho,sec:specgap,sec:Sigma,sec:three-term,sec:task-null,sec:bilevel}
treats $\rcontact$, $\specgap$, and the contact force $f = (f_{n},
f_{t})$ as quantities that vary continuously along admissible
trajectories, with the rigid-body limit recovered only as a
singular limit. The present subsection makes the modeling regime
explicit: it names the three formulations of contact dynamics in
current use, identifies the bridge between them, and states which
formulation the structural results of §2 are written against.

\paragraph{Three formulations of contact dynamics.}
Contact dynamics admits three formulations that are mathematically
equivalent in appropriate limits, each preferred in a different
computational setting:
\begin{enumerate}[label=(\arabic*),leftmargin=*]
  \item \emph{Set-valued continuous-time formulation.} The
        contact-coupled dynamics is written as a measure-valued
        differential inclusion
        \[
          \dot v \;\in\; F(t, q, v) - \mathcal{N}_{K}(v),
        \]
        where $\mathcal{N}_{K}$ is the normal cone to the admissible
        velocity set induced by the unilateral contact constraints
        and the friction cone $\frictioncone$. This is the
        foundational treatment of \citet{Stewart2000} and the
        modern textbook formulation of \citet{Brogliato2016}.
  \item \emph{Contact-implicit trajectory optimization (CITO).}
        Time-discretize the system over a horizon $\{t_{k}\}$ and
        include the discrete-time complementarity
        \[
          0 \;\le\; \mathbf{f}_{k} \;\perp\; \phi_{k} \;\ge\; 0
        \]
        as a hard constraint in the trajectory nonlinear program,
        yielding a Mathematical Program with Complementarity
        Constraints (MPCC). The foundational time-stepping LCP is
        \citet{StewartTrinkle1996}; the canonical trajectory-level
        formulation through contact is
        \citet{PosaCantuTedrake2014}.
  \item \emph{Per-timestep convex relaxation.} Replace the
        discrete-time complementarity at each step with a strictly
        convex regularization, yielding a per-timestep convex
        QP/SOCP whose solution converges to the LCP solution as the
        regularization parameter goes to zero. The foundational
        optimization-based simulation is \citet{Anitescu2006}; the
        convex contact model used by current dynamics simulators is
        \citet{Todorov2014}. (A penalty/spring--damper formulation
        $f_{n} = K\, \phi_{+} + B\, \dot\phi$ is a different
        smoothing regime and is not used in this paper.)
\end{enumerate}

\paragraph{Bridge: set-valued and MPCC are dual formulations.}
The set-valued formulation~(1) and the MPCC
formulation~(2) are connected by an explicit identity. The
Signorini complementarity admits the normal-cone form
\begin{equation}
\label{eq:signorini-nc}
  0 \;\le\; f_{n} \;\perp\; \phi \;\ge\; 0
  \quad\Longleftrightarrow\quad
  -\, f_{n} \;\in\; \mathcal{N}_{\mathbb{R}_{+}}(\phi),
\end{equation}
and the friction subdifferential inclusion \eqref{eq:friction-di}
of \cref{sec:bilevel} writes the tangential force as
$f_{t} \in -\mu\, f_{n}\, \partial \|v_{t}\|$. Together,
\eqref{eq:signorini-nc} and \eqref{eq:friction-di} are the
unilateral and tangential components of the set-valued formulation
(1). Time-discretizing them with an implicit Euler scheme yields
per-step linear complementarity problems whose
trajectory-horizon composition is the MPCC of formulation~(2).
Convex regularization of those per-step LCPs in turn yields
formulation~(3). The three formulations are therefore a
\emph{continuous-time} $\to$
\emph{discrete-time-trajectory} $\to$
\emph{discrete-time-per-step-convex-regularized} chain, each step
approximating the previous one. The friction subdifferential
inclusion already written in \cref{sec:bilevel} is itself a
member of formulation~(1); the present subsection makes the
chain to (2) and (3) explicit so that the modeling regime of the
subsequent subsections is unambiguous.

\paragraph{The framework's modeling regime.}
The framework of \cref{sec:rho,sec:specgap,sec:Sigma,sec:three-term,sec:task-null,sec:bilevel}
is consistent with formulation~(3): per-timestep convex
relaxation of the rigid-body LCP. Under this regime, the contact
force $f_{n}$ is a continuous function of state and time
(no impulsive Dirac contributions inside the per-step solve), so
$\rcontact$ is continuous along admissible trajectories and the
hypothesis of the two-crossing lemma (\cref{lem:two-crossing}) is
satisfied. The contact-graph edge weights $w_{ij}$ of
\cref{sec:specgap} are finite---they are the regularized contact
stiffnesses produced by the convex per-step solve. The bang arcs
of \cref{sec:three-term} are high-rate but not literally
impulsive: the spectral-kick term $u_{s}$ is a high-impulse but
absolutely continuous control input on a small interval
surrounding $t_{1}$ and $t_{4}$, not a Dirac mass. This is the
regime simulated by MuJoCo \citep{Todorov2014} and physically
realized by manipulators with finite finger compliance, including
the Franka platform of the experimental staircase.

\paragraph{Why CITO/MPCC is not required here.}
CITO/MPCC trajectory optimization (formulation~(2)) solves a
problem the present framework does not pose: it discovers the
contact mode---which contacts are active in which time
intervals---jointly with the trajectory, by treating the
complementarity as an optimization constraint. The present
framework instead \emph{pre-specifies} the contact mode through
the phase decomposition of \cref{sec:phase-cycle} (forward
reference) and the plan generator of \cref{sec:plan-training}
(forward reference): phase~1 and phase~5 are open-contact
intervals (no robot--object contact); phases~2 and~4 are
contact-mode transitions localized at $t_{1}$ and $t_{4}$;
phase~3 is a closed-contact interval. Given this
pre-specification, the trajectory-level optimization of
\cref{def:bilevel} need not search over contact modes---those are
determined by the plan---and what remains is the per-timestep
convex contact solve of formulation~(3), which is exactly what
MuJoCo and similar simulators provide as a primitive. CITO is
therefore a complementary tool that solves a problem the present
framework chooses to side-step by leveraging the structural
pre-specification of phase order; the two approaches are not in
conflict but target different decompositions of the manipulation
problem. \emph{The contact-mode pre-specification of the phase
decomposition is the structural reason this paper does not
require an outer-loop MPCC formulation.}

\paragraph{Limit stability.}
The structural results of §2---the five-phase decomposition
(\cref{lem:two-crossing}), the three-term decomposition
(\cref{thm:three-term}), the bi-level scaffold
(\cref{def:bilevel}), and the spectral-gap interpretation
(\cref{lem:specgap-interp})---are stable as the convex
regularization parameter goes to zero. The rigid-body LCP limit
is recovered as a singular limit in which $u_{s}$ becomes a Dirac
mass and $\rcontact$ acquires one-sided continuity at $t_{1}$
and $t_{4}$; the structural conclusions persist under
measure-theoretic interpretation. The choice of regime above is
therefore modeling convenience, not a structural restriction of
the framework. Detailed limit analysis is out of scope for this
paper.

% ---------------------------------------------------------------
\subsection{Causal observability}
\label{sec:causal-obs}

The framework of \cref{sec:rho,sec:specgap,sec:Sigma,sec:three-term,sec:task-null,sec:bilevel}
treats the contact order parameter $\rcontact$, the contact graph
$G_{c}$, and the object pose as quantities known to the controller
at every instant. In a real manipulation pipeline these
quantities are not given but estimated, and the estimates flow
through sensors with limited field of view (FOV), latency, and
noise. The present subsection adds the perception-side companion
of the contact-side framework: where $\rcontact$ is the central
\emph{interface variable} of \cref{sec:rho}, the
\emph{observation map} $\obsmap(\cdot)$ introduced below is the
central perception-side interface, and the notion of
\emph{causal observability} formalized in
\cref{def:causal-obs} states the temporal precondition under
which the perception subsystem can deliver state estimates in
time for the controller to act on them.

\begin{definition}[Observation map]\label{def:obsmap}
For a manipulator state $s \in \mathcal{S}$ (joint configuration,
velocity, and any object pose information present in the scene)
and a sensor reading $y$ in some sensor space $\mathcal{Y}$, the
\emph{observation map}
\[
  \obsmap : \mathcal{S} \to \mathcal{Y}
\]
is the deterministic noise-free mapping from manipulator state to
sensor reading. A real sensor produces $y_{t} = \obsmap(s_{t}) +
\nu_{t}$ where $\nu_{t}$ is bounded measurement noise; we suppress
$\nu$ for the present subsection and add it back when analyzing
specific modalities in \cref{sec:perception}.
\end{definition}

\paragraph{Cognitive boundary.}
Not every scene quantity produces a sensor reading from which it
can be uniquely recovered. For a wrist-mounted camera, only
objects within the camera's FOV cone at the current configuration
project onto the image plane at all; for proprioception, only
joints with non-degenerate Jacobian rows are individually
distinguishable; for haptics, only the contact patch in immediate
contact with the fingertip is sensed. The observable subset of
the scene is configuration-dependent.

\begin{definition}[Cognitive boundary]\label{def:cognitive-boundary}
The \emph{cognitive boundary} of the observation map $\obsmap$
at configuration $\jointvec$ is the subset
\[
  D(\jointvec)
  \;\subseteq\; \mathbb{R}^{3} \times SE(3)
\]
of object positions and orientations whose images under
$\obsmap$ are distinguishable from one another given $\jointvec$.
For a wrist-mounted camera, $D(\jointvec)$ is the projection
into world coordinates of the camera's FOV cone at the current
configuration.
\end{definition}

The cognitive boundary $D(\jointvec)$ is a static notion: at a
fixed configuration, it tells us which scene quantities are even
in principle recoverable. Manipulation control, however, runs
against task deadlines determined by the phase structure of
\cref{sec:Sigma}, and a scene quantity that is recoverable in
principle but only after the relevant deadline has passed is, for
control purposes, indistinguishable from one that is not
recoverable at all. The temporal version of distinguishability is
the central concept of this subsection.

\begin{definition}[Causal observability]\label{def:causal-obs}
A scene quantity $r$ (e.g.\ an object's pose) is \emph{causally
observable} at time $t$ along a trajectory $\jointvec(\cdot)$ if
there exists a time $t' \le t + \tau_{\mathrm{deadline}}$ such
that
\begin{enumerate}[label=(\roman*),nosep]
  \item $r$ lies in the cognitive boundary $D(\jointvec(t''))$ for
        some $t'' \in [t, t']$ (the sensor can capture it);
  \item the latency budget
        $\tau_{\mathrm{sensor}} + \tau_{\mathrm{compute}}
        \le t' - t''$
        is respected (the estimate reaches the controller before
        the deadline).
\end{enumerate}
A scene quantity that fails (i) is \emph{spatially invisible}; one
that fails only (ii) is \emph{visible but late}. \emph{Causally
invisible} means it fails at least one of (i), (ii).
\end{definition}

\paragraph{Relation to standard frameworks.}
\Cref{def:causal-obs} is partial observability with explicit
deadline awareness. It is the manipulation specialization of the
partially observable Markov decision process (POMDP) framework
\citep{KaelblingLittmanCassandra1998}, in which the state is
hidden but partially recoverable through observations; the
manipulation specialization adds the requirement that the
observation be obtained before a \emph{task-relevant} deadline.
The relevant deadline in our setting is the next
$\swsurf$-crossing time of \cref{sec:Sigma}, or any other
phase-determined event time identified by the phase decomposition
of \cref{sec:three-term}. Standard POMDP analysis quantifies
\emph{what} is recoverable from a given observation history;
causal observability adds the \emph{when}.

\paragraph{Observation--action coupling.}
In manipulation, action moves the FOV: the wrist-mounted camera
rides the end-effector, so the cognitive boundary
$D(\jointvec)$ depends on the trajectory $\jointvec(\cdot)$.
This makes observation closed-loop: what is causally observable
at time $t$ depends on actions taken before $t$. Formally, the
observation history is
$y_{0:t} = \obsmap(s_{0:t})$
where $s_{0:t}$ depends on the control history $u_{0:t-1}$
through the dynamics of \cref{sec:task-null}. A robot that has
not yet looked at the object cannot have observed it, and a
robot that has not yet observed the object cannot plan a look
that is sure to bring it into the FOV. This is the manipulation
analog of \emph{active perception} \citep{Bajcsy1988}, in which
the sensor's pose is itself a control variable.

\paragraph{Information timeliness.}
A causally observable scene quantity has positive control value
if and only if the latency budget closes the loop in time:
\begin{equation}
\label{eq:timeliness}
  \tau_{\mathrm{sensor}} + \tau_{\mathrm{compute}}
  \;<\; \tau_{\mathrm{deadline}},
\end{equation}
where the three latencies are tabulated in \cref{tab:notation}:
$\tau_{\mathrm{sensor}}$ is the time from photon arrival to
sensor read-out, $\tau_{\mathrm{compute}}$ is the time from
sensor read-out to a state estimate available at the controller,
and $\tau_{\mathrm{deadline}}$ is the time until the next
phase-determined event the estimate is needed for. The strict
inequality \eqref{eq:timeliness} is the temporal analog of the
spectral-gap margin condition $\specgap \ge \varepsilon$ of
\cref{sec:specgap}: the controller maintains a quantitative
force margin during transport and a quantitative latency margin
during perception, and both margins are state-dependent
quantities the controller actively manages.

\paragraph{Forward references.}
\Cref{sec:perception} specializes the observation map to three
sensory modalities (vision, proprioception, haptics), each with
its own cognitive boundary $D(\jointvec)$ and its own latency
profile, and treats the stochastic extension in which $D$ is
replaced by a probability-of-detection field.
\Cref{sec:plan-training} uses causal observability as a
feasibility precondition on plans: a plan that requires the
controller to act on a causally invisible scene quantity at any
phase boundary is rejected before being passed to the
controller.

% ---------------------------------------------------------------
\subsection{Notation}
\label{sec:notation}

Throughout the paper we use the symbols collected in
\cref{tab:notation}. Where a symbol is overloaded across the
robotics literature, we adopt the convention closest to the
manipulation context (Khatib--Nakamura for task/null-space,
Brogliato for nonsmooth contact, Lipman et al.\ for the trained
velocity field). Two symbol overloads internal to this paper are
called out below the table.

\begin{table}[h]
\centering
\small
\begin{tabular}{@{}lll@{}}
\toprule
Symbol & Meaning & First introduced \\
\midrule
\multicolumn{3}{@{}l}{\textit{Order parameter and switching}}\\
$\rcontact$                & contact order parameter
                           & \S2.1 \\
$\swsurf = \{\rcontact = 1\}$ & switching surface
                           & \S2.3 \\
$\specgap$                 & spectral gap of contact graph Laplacian
                           & \S2.2 \\
$\specgapeff$              & effective spectral gap with wind correction
                           & \S7.2 \\
$\varepsilon$              & spectral gap threshold ($\specgap \ge \varepsilon$ on contact)$^{\dagger}$
                           & \S3.1 \\
\midrule
\multicolumn{3}{@{}l}{\textit{Phase structure}}\\
$\phaseop{0},\dots,\phaseop{4}$
                           & phase operators (pick / load / use / unload / place)
                           & \S3.2 \\
$\varphi$                  & phase indicator (one-hot in $\mathbb{R}^{5}$)
                           & \S6.3 \\
$\Greflect$                & reflection group $\{z \mapsto -z,\ g \mapsto -g,\ F \mapsto -F\}$
                           & \S3.3 \\
\midrule
\multicolumn{3}{@{}l}{\textit{Robot kinematics and dynamics}}\\
$\jointvec, \dot{\jointvec}, \ddot{\jointvec} \in \mathbb{R}^{n}$
                           & joint configuration, velocity, acceleration ($n = 7$ for Franka)
                           & \S2.5 \\
$J(\jointvec)$             & task Jacobian
                           & \S2.5 \\
$J^{+}$                    & Moore--Penrose pseudo-inverse
                           & \S2.5 \\
$M(\jointvec), C(\jointvec, \dot{\jointvec}), g(\jointvec)$
                           & inertia, Coriolis, gravity terms
                           & \S4.2 \\
$\InvDyn$                  & inverse dynamics operator
                           & \S6.3 \\
$u$                        & joint torque control input
                           & \S2.4 \\
\midrule
\multicolumn{3}{@{}l}{\textit{Task-space / null-space}}\\
$\TS$                      & task-space subspace (range of $J^{\top}$)
                           & \S2.5 \\
$\NS$                      & null-space subspace (kernel of $J$)
                           & \S2.5 \\
$[\TS, \NS]$               & dexterity commutator
                           & \S4.3 \\
\midrule
\multicolumn{3}{@{}l}{\textit{Contact and forces}}\\
$f = (f_{n}, f_{t})$       & contact force (normal, tangential)
                           & \S4.2 \\
$\frictioncone$            & friction cone
                           & \S4.2 \\
$\mu$                      & friction coefficient
                           & \S4.2 \\
$\Fwind, \Fmax$            & adversarial wind force, ceiling
                           & \S7.2 \\
\midrule
\multicolumn{3}{@{}l}{\textit{Perception}}\\
$\obsmap(\cdot)$           & observation map
                           & \S2.7 \\
$\Wreach$                  & reachable workspace
                           & \S2.7 \\
$D(\jointvec)$             & camera FOV at configuration $\jointvec$
                           & \S5.1 \\
$\tau_{\mathrm{sensor}}, \tau_{\mathrm{compute}}, \tau_{\mathrm{deadline}}$
                           & timeliness latencies
                           & \S5.3 \\
\midrule
\multicolumn{3}{@{}l}{\textit{Plan and training}}\\
$P, P^{*}$                 & plan, optimal plan
                           & \S6.1 \\
$\veltheta$                & CFM-trained velocity field
                           & \S6.2 \\
$\polphi$                  & RLMD-trained policy
                           & \S6.4 \\
$\actpsi$                  & actuator residual net
                           & \S6.4 / \S8.2 \\
$\mmdloss(\theta)$         & MMD loss
                           & \S6.2 \\
\midrule
\multicolumn{3}{@{}l}{\textit{RLMD (risk lock)}}\\
$\risklock$                & risk lock threshold
                           & \S6.4 \\
$\softrisk$                & soft-surrogate failure rate
                           & \S6.4 \\
$\actrough$                & action roughness penalty
                           & \S6.4 \\
$\beta$                    & sigmoid temperature for $\softrisk$
                           & \S6.4 \\
$\curriculumeps$           & curriculum perturbation at iteration $k$$^{\dagger}$
                           & \S6.4 \\
\midrule
\multicolumn{3}{@{}l}{\textit{Curriculum}}\\
$L$                        & stack layers
                           & \S7.1 \\
$N_{\mathrm{blocks}}$      & number of blocks$^{\ddagger}$
                           & \S7.1 \\
$K_{\mathrm{type}}$        & number of block types
                           & \S7.1 \\
\bottomrule
\end{tabular}
\caption{Notation. Subsection references point to the first
introduction of the symbol.
$^{\dagger}$~The two symbols $\varepsilon$ (spectral gap threshold,
\S3.1) and $\curriculumeps$ (curriculum perturbation at iteration $k$,
\S6.4) share the Greek letter $\varepsilon$ but denote distinct
quantities; we always disambiguate by the subscript $k$.
$^{\ddagger}$~The base letter $N$ is also used in \S6.4 for the
RLMD rollout count $N_{\mathrm{rollout}}$; both occurrences carry
explicit subscripts in the body to avoid ambiguity.}
\label{tab:notation}
\end{table}

\begin{remark}[Symbol overloads]
\label{rem:overloads}
Two notational collisions are unavoidable in a manipulation paper
of this scope and are flagged here for the reader.
\emph{(i)} The Greek letter $\varepsilon$ appears as the spectral
gap threshold $\specgap \ge \varepsilon$ in \S3.1 and as the
indexed curriculum perturbation $\curriculumeps$ in \S6.4. The
subscript $k$ disambiguates the latter and is always present.
\emph{(ii)} The base letter $N$ appears as $N_{\mathrm{blocks}}$
(the number of blocks in the stacking task, \S7.1) and as
$N_{\mathrm{rollout}}$ (the RLMD rollout count, \S6.4); both
occurrences carry explicit subscripts and the bare symbol $N$ is
not used in the body.
\end{remark}
